\section{总结与延伸}

\subsection{核心思想图谱}

胡渊鸣的思维方式可以用几个核心信念来概括：

\textbf{关于世界观}：这个世界很可能是一个模拟——普朗克常数是浮点精度，光速是计算资源限制。作为被模拟的一员，我们无法从内部证伪这个假说，但这种视角深刻影响了他对仿真、AI 和真实世界关系的理解。

\textbf{关于技术}：物理定律是准的，但边界条件不是。未来的仿真将以 learning 为主、物理先验为辅。3D inductive bias 是连接 simulation 和 neural network 的桥梁。不要低估技术前进的速度。

\textbf{关于创业}：选对方向比努力重要十倍。技术必须能商业化才有持续性。开源可以是战略武器，但需要与商业化互补。创业公司要找中等规模、非共识的市场。

\textbf{关于个人成长}：99\%是运气（包括能努力本身）。只做自己 enjoy 的事情。每三个月把过去的自己否定掉。多花时间想 what 和 why，而不只是 how。不要追求所有人喜欢你。

\subsection{对不同读者的启示}

\textbf{对科研工作者}：先花95\%的时间想什么问题是重要的问题，再花5\%的时间去解决（引自爱因斯坦）。发论文之前，先问自己"这个工作重要吗？"——这也是胡渊鸣面试每个 Meshy 求职者时会问的问题。不要把自己关在办公室里，把门开着，多听其他领域的人在做什么。阶段一该卷就卷，但有了3-4篇顶会论文后，必须切换到阶段二思维——不是发更多论文，而是做更有 impact 的工作。

\textbf{对创业者}：商业化是第一性原理——技术不能商业化就是"烟花秀"。性格决定商业模式——内向型创始人做标准化消费级产品（subscription + API），避免做需要大量客户关系维护的外包。Pivot 不可怕，可怕的是没有 pivot 的勇气。面对事实、尊重数据，不要被"天道酬勤"的鸡汤欺骗。好生意的三要素：市场大、非共识（timing）、有竞争优势。选方向比努力重要十倍。

\textbf{对管理者}：好的 leader 不一定被所有人喜欢，但应该被所有人服气。新任管理者最大的坎是敢于对昔日同级提出更高标准。公司文化为"赢"服务——内耗、甩锅、追求和谐都是"会输的文化"。CEO 最重要的事：40\%招聘、30\%思考未来战略、30\%参与日常业务。保持 hands-on——不写代码的 CEO 容易"飘"。

\textbf{对所有人}：
\begin{itemize}
\item 找到自己真正 enjoy 的事情，越早越好
\item 99\%是运气，保持谦逊
\item 每三个月审视一次自己——如果不觉得三个月前的自己是傻逼，说明没有成长
\item 不要追求所有人都喜欢你——把精力放在做正确的事情上
\item 真正的失败不是项目做砸了，而是不敢再尝试了
\item 要有勇气——"希望所有看到这个 Podcast 的人都能具备勇气"
\end{itemize}

\subsection{胡渊鸣语录精选}

以下是访谈中最具启发性的原话摘录：

\begin{table}[H]
\centering
\caption{胡渊鸣核心语录}
\begin{tabular}{p{3.5cm}p{10cm}}
\toprule
\textbf{主题} & \textbf{原话} \\
\midrule
自我定义 & "我是一个被自己的兴趣和使命感驱动着去做事的人。" \\
\midrule
运气 & "99\%都是运气。甚至你能努力也是运气的一部分。" \\
\midrule
创业成长 & "每三到六个月就要把自己开除一次。最好就是觉得三个月前的自己是傻逼。" \\
\midrule
CEO 觉悟 & "如果你工作的目标是让所有人都喜欢你，那你最终会被所有人讨厌。" \\
\midrule
开源价值 & "开源确实不挣钱，但它能帮你吸引最好的人才和客户。" \\
\midrule
努力与回报 & "世界上绝大多数努力没有任何回报。天道酬勤完全是骗人的。" \\
\midrule
失败观 & "只要我还活着，我就没有失败。真正的失败是不愿意再去创新。" \\
\midrule
创业锻炼 & "和平年代没有什么比创业更能锻炼一个人了。" \\
\midrule
模拟假说 & "也许整个世界就是被一个更高级文明闲着无聊写的程序来模拟的。" \\
\midrule
AI 与物理 & "Simulator 里面需要更多 AI，AI 里面也需要更多 3D inductive bias。" \\
\bottomrule
\end{tabular}
\end{table}

\subsection{拓展阅读}

\begin{itemize}
\item \textbf{太极编程语言}：\url{https://github.com/taichi-dev/taichi}——GitHub Stars 超过 NVIDIA Warp、Triton、Mojo
\item \textbf{Meshy AI}：\url{https://www.meshy.ai/}——当前市场占有率第一的3D生成平台
\item \textbf{胡渊鸣知乎文章}：关于"如果重读PhD"的反思，以及"不要只做 SIGGRAPH 里的小鱼"的论述
\item \textbf{Andy Grove}：《只有偏执狂才能生存》（Only the Paranoid Survive）——Intel 从 DRAM 转型 CPU 的决策过程
\item \textbf{Richard Hamming}：You and Your Research——关于"把门开着"的经典演讲
\item \textbf{Sergey Levine}：The Spork of AGI——关于 simulation 在 robotics 中的角色与局限
\item \textbf{Andrej Karpathy}：关于"十倍收益往往只有两倍难"的观点
\item \textbf{Games201}：胡渊鸣在 COVID 期间开设的物理仿真课程（中文），他计划未来推出"高清重制版"
\item \textbf{DiffTaichi (ICLR 2020)}：太极可微编程的标志性论文，将 simulation 和 robotics 结合
\item \textbf{ChainQueen}：太极在可微仿真领域的早期工作
\item \textbf{QuantTaichi}：实现单GPU 10亿粒子仿真的量化优化工作
\item \textbf{雷军}：《小米创业思考》——胡渊鸣推荐的创业参考，"飞猪理论"与方向选择的重要性
\end{itemize}

%% --- Fin del contenido --- %%

\end{document}
